% GENERATED FILE. Edit canonical lessons, then run npm run generate:books.
% canonical-source-hash: fnv1a64:6ef24c6261640f44
% canonical-lessons: MR-C424-vikri, MR-C424-pherivala, MR-C424-ghauk, MR-C424-dalal, MR-C424-bil

\chapter{Sellers and Hawkers}
\label{ch:mr-sellers-and-hawkers}
\hlchaptermodality{424}

\begin{chapteropening}
\textbf{By the end of this chapter:} \emph{I can talk about sales, hawkers, wholesale buying, brokers and bills.}

Complete the last lesson of chapter 424: I can talk about sales, hawkers, wholesale buying, brokers and bills.
\end{chapteropening}

\section[\texorpdfstring{\mr{विक्री} (vikrī)}{विक्री (vikrī)}]{\mr{विक्री} (vikrī) — a sale; selling}

\label{lesson:MR-C424-vikri}



\begin{quote}
Before the new one: say the Marathi for to bargain, then the Marathi for to buy.
\end{quote}

\subsection*{You'll want to know: \mr{विक्री}}
\textbf{\mr{विक्री}} — \emph{vikrī} — \textquotedblleft{}a sale; selling\textquotedblright{}.

\begin{quote}
\mr{या} \mr{महिन्यात} \mr{दुकानाची} \mr{विक्री} \mr{वाढली}.
\end{quote}

\begin{quote}
\emph{yā mahinyāt dukānācī vikrī vāḍhlī.}
\end{quote}

\begin{quote}
This month the shop’s sales went up.
\end{quote}

Feminine.

\mr{विक्रीला}, for sale: \mr{घर} \mr{विक्रीला} \mr{आहे}, the house is for sale.

\begin{cousinweb}
\textbf{\mr{विक्री}} (\emph{vikrī}) is inherited: CDIAL lists Marathi \emph{vikrī}, \textquotedblleft{}sale\textquotedblright{}, under Sanskrit \emph{vikraya}, \textquotedblleft{}sale\textquotedblright{} (CDIAL 11638).
\end{cousinweb}

\begin{grammarlens}[title={What you've built}]
One more word for getting around.
\end{grammarlens}

\subsection*{Your turn}
\begin{itemize}
\raggedright
  \item \emph{Say it:} \emph{vikrī}
  \item \emph{Say it:} \emph{vikrī}, once more
  \item \emph{Say it:} \emph{yā mahinyāt dukānācī vikrī vāḍhlī.} — This month the shop’s sales went up.
  \item \emph{Recall:} say the Marathi for to bargain, then the Marathi for to buy, then say \emph{vikrī} again
\end{itemize}

\begin{tcolorbox}[breakable,colback=teal!4,colframe=teal!35!black,title={Before you move on}]
What does \mr{विक्री} mean? (\textquotedblleft{}A sale; selling\textquotedblright{}.) Say it once more.
\end{tcolorbox}
\section[\texorpdfstring{\mr{फेरीवाला} (pherīvālā)}{फेरीवाला (pherīvālā)}]{\mr{फेरीवाला} (pherīvālā) — a hawker}

\label{lesson:MR-C424-pherivala}



\begin{quote}
Before the new one: say the Marathi for to buy, then the Marathi for a sale.
\end{quote}

\subsection*{You'll want to know: \mr{फेरीवाला}}
\textbf{\mr{फेरीवाला}} — \emph{pherīvālā} — \textquotedblleft{}a hawker\textquotedblright{}.

\begin{quote}
\mr{फेरीवाला} \mr{गल्लीत} \mr{भाजी} \mr{विकतो}.
\end{quote}

\begin{quote}
\emph{pherīvālā gallīt bhājī vikto.}
\end{quote}

\begin{quote}
The hawker sells vegetables in the lane.
\end{quote}

Masculine.

\mr{फेरी}, a round, and -\mr{वाला}, the one who does it.

\begin{cousinweb}
\textbf{\mr{फेरीवाला}} (\emph{pherīvālā}) is built from \textbf{\mr{फेरी}}, \textquotedblleft{}a round\textquotedblright{}, from the family of \emph{*phirati}, \textquotedblleft{}turns\textquotedblright{} (CDIAL 9078), and -\mr{वाला}, \textquotedblleft{}the one who\textquotedblright{}.
\end{cousinweb}

\begin{grammarlens}[title={What you've built}]
One more word for getting around.
\end{grammarlens}

\subsection*{Your turn}
\begin{itemize}
\raggedright
  \item \emph{Say it:} \emph{pherīvālā}
  \item \emph{Say it:} \emph{pherīvālā}, once more
  \item \emph{Say it:} \emph{pherīvālā gallīt bhājī vikto.} — The hawker sells vegetables in the lane.
  \item \emph{Recall:} say the Marathi for to buy, then the Marathi for a sale, then say \emph{pherīvālā} again
\end{itemize}

\begin{tcolorbox}[breakable,colback=teal!4,colframe=teal!35!black,title={Before you move on}]
What does \mr{फेरीवाला} mean? (\textquotedblleft{}A hawker\textquotedblright{}.) Say it once more.
\end{tcolorbox}
\section[\texorpdfstring{\mr{घाऊक} (ghāūk)}{घाऊक (ghāūk)}]{\mr{घाऊक} (ghāūk) — wholesale}

\label{lesson:MR-C424-ghauk}



\begin{quote}
Before the new one: say the Marathi for a sale, then the Marathi for a hawker.
\end{quote}

\subsection*{You'll want to know: \mr{घाऊक}}
\textbf{\mr{घाऊक}} — \emph{ghāūk} — \textquotedblleft{}wholesale\textquotedblright{}.

\begin{quote}
\mr{तो} \mr{घाऊक} \mr{बाजारातून} \mr{कांदे} \mr{आणतो}.
\end{quote}

\begin{quote}
\emph{to ghāūk bājārātūn kānde āṇto.}
\end{quote}

\begin{quote}
He brings onions from the wholesale market.
\end{quote}

It does not change: \mr{घाऊक} \mr{बाजार}, \mr{घाऊक} \mr{खरेदी}.

\begin{cousinweb}
\textbf{\mr{घाऊक}} (\emph{ghāūk}) is a Marathi word for buying and selling in bulk; its origin is not settled.
\end{cousinweb}

\begin{grammarlens}[title={What you've built}]
One more word for getting around.
\end{grammarlens}

\subsection*{Your turn}
\begin{itemize}
\raggedright
  \item \emph{Say it:} \emph{ghāūk}
  \item \emph{Say it:} \emph{ghāūk}, once more
  \item \emph{Say it:} \emph{to ghāūk bājārātūn kānde āṇto.} — He brings onions from the wholesale market.
  \item \emph{Recall:} say the Marathi for a sale, then the Marathi for a hawker, then say \emph{ghāūk} again
\end{itemize}

\begin{tcolorbox}[breakable,colback=teal!4,colframe=teal!35!black,title={Before you move on}]
What does \mr{घाऊक} mean? (\textquotedblleft{}Wholesale\textquotedblright{}.) Say it once more.
\end{tcolorbox}
\section[\texorpdfstring{\mr{दलाल} (dalāl)}{दलाल (dalāl)}]{\mr{दलाल} (dalāl) — a broker; an agent}

\label{lesson:MR-C424-dalal}



\begin{quote}
Before the new one: say the Marathi for a hawker, then the Marathi for wholesale.
\end{quote}

\subsection*{You'll want to know: \mr{दलाल}}
\textbf{\mr{दलाल}} — \emph{dalāl} — \textquotedblleft{}a broker; an agent\textquotedblright{}.

\begin{quote}
\mr{दलालाने} \mr{आम्हाला} \mr{भाड्याचं} \mr{घर} \mr{दाखवलं}.
\end{quote}

\begin{quote}
\emph{dalālāne āmhālā bhāḍyācaṁ ghar dākhavlaṁ.}
\end{quote}

\begin{quote}
The broker showed us a house for rent.
\end{quote}

Masculine.

\mr{दलाली} is his fee.

\begin{cousinweb}
\textbf{\mr{दलाल}} (\emph{dalāl}) is a loan through Persian, usually traced to Arabic \emph{dallāl}, \textquotedblleft{}a broker, a guide\textquotedblright{}.
\end{cousinweb}

\begin{grammarlens}[title={What you've built}]
One more word for getting around.
\end{grammarlens}

\subsection*{Your turn}
\begin{itemize}
\raggedright
  \item \emph{Say it:} \emph{dalāl}
  \item \emph{Say it:} \emph{dalāl}, once more
  \item \emph{Say it:} \emph{dalālāne āmhālā bhāḍyācaṁ ghar dākhavlaṁ.} — The broker showed us a house for rent.
  \item \emph{Recall:} say the Marathi for a hawker, then the Marathi for wholesale, then say \emph{dalāl} again
\end{itemize}

\begin{tcolorbox}[breakable,colback=teal!4,colframe=teal!35!black,title={Before you move on}]
What does \mr{दलाल} mean? (\textquotedblleft{}A broker; an agent\textquotedblright{}.) Say it once more.
\end{tcolorbox}
\section[\texorpdfstring{\mr{बिल} (bil)}{बिल (bil)}]{\mr{बिल} (bil) — a bill}

\label{lesson:MR-C424-bil}



\begin{quote}
Before the new one: say the Marathi for wholesale, then the Marathi for a broker.
\end{quote}

\subsection*{You'll want to know: \mr{बिल}}
\textbf{\mr{बिल}} — \emph{bil} — \textquotedblleft{}a bill\textquotedblright{}.

\begin{quote}
\mr{जेवणानंतर} \mr{मी} \mr{बिल} \mr{भरलं}.
\end{quote}

\begin{quote}
\emph{jevṇānantar mī bil bharlaṁ.}
\end{quote}

\begin{quote}
After the meal I paid the bill.
\end{quote}

Neuter.

\mr{बिल} \mr{भरणे}, to pay a bill.

\begin{cousinweb}
\textbf{\mr{बिल}} (\emph{bil}) is a loan from English \emph{bill}.
\end{cousinweb}

\begin{grammarlens}[title={What you've built}]
One more word for getting around.
\end{grammarlens}

\subsection*{Your turn}
\begin{itemize}
\raggedright
  \item \emph{Say it:} \emph{bil}
  \item \emph{Say it:} \emph{bil}, once more
  \item \emph{Say it:} \emph{jevṇānantar mī bil bharlaṁ.} — After the meal I paid the bill.
  \item \emph{Recall:} say the Marathi for wholesale, then the Marathi for a broker, then say \emph{bil} again
\end{itemize}

\begin{tcolorbox}[breakable,colback=teal!4,colframe=teal!35!black,title={Before you move on}]
What does \mr{बिल} mean? (\textquotedblleft{}A bill\textquotedblright{}.) Say it once more.
\end{tcolorbox}
